% =====================================================================
%  3.5 Eksponential- og potensfunksjoner — ekstra utfordring
% =====================================================================

\begin{infoboks}{Velg metode selv}
Arket bygger på 3.5 Eksponential- og potensfunksjoner. Vis regning og begrunnelse.
Tallene i U1 er tenkte. Kleibers lov i U6 er en virkelig modell, laget ut fra målinger. Bruk eksakte svar når det går.
\end{infoboks}

\nivaahode{2}

\oppg{U1}{2}{\uh}
Tenkt eksempel: En medisin har halveringstid 6 timer. Startdosen er 400 mg. Mengden etter $t$ timer er
\[ f(t) = 400 \cdot \left(\tfrac{1}{2}\right)^{\tfrac{t}{6}} \]
\begin{tasks}(2)
  \task! Forklar hvorfor funksjonen passer med halveringstiden 6 timer.
  \task Hvor mye er igjen etter ett døgn?
  \task Hvor mye er igjen etter 18 timer?
\end{tasks}

\oppg{U2}{2}{\uh}
Løs likningene. Skriv begge sider med samme grunntall først.
\begin{tasks}(3)
  \task $5^x = 125$
  \task $3 \cdot 2^x = 96$
  \task $\left(\dfrac{1}{2}\right)^x = 16$
\end{tasks}

\oppg{U3}{2}{\prog}
Programmet finner hvor mange hele år et beløp trenger for å bli doblet ved $p\,\%$ rente. Det skriver også ut $\frac{72}{p}$ (72-regelen).
\begin{kode}
for p in range(1, 11):
    belop = 1.0
    ar = 0
    while belop < 2:
        belop = belop * (1 + p / 100)
        ar = ar + 1
    print(p, ar, 72 / p)
\end{kode}
\begin{tasks}(1)
  \task Forklar hva som skjer i den indre løkken. Hvorfor starter \texttt{belop} på 1,0?
  \task Kjør programmet. Hva skriver det ut for $p = 4$ og $p = 7$?
  \task Hvor godt stemmer 72-regelen? Regelen er en tommelfingerregel og ikke eksakt. Kommenter avvikene.
\end{tasks}

\nivaahode{3}

\oppg{U4}{3}{\uh}
Bevis.
\begin{tasks}(1)
  \task Vis at hvis $\dfrac{f(x+1)}{f(x)} = b$ for alle heltall $x$, er $f(n) = f(0) \cdot b^n$ for alle heltall $n \ge 0$.
  \task Anta at $\dfrac{g(2x)}{g(x)} = r$ for $x = 1, 2, 4, 8, \ldots$ Vis at $g(2^n) = g(1) \cdot r^n$. Hva betyr det når $r = 2^k$?
  \task Hva betyr de to resultatene for hvordan en tabell ser ut?
\end{tasks}

\oppg{U5}{3}{\uh}
Er det mulig at en funksjon både kan skrives som $f(x) = a \cdot b^x$ og som $f(x) = c \cdot x^k$ for alle $x > 0$?
(Vi tillater $b = 1$ og $k = 0$.) Undersøk, og finn alle slike funksjoner. Begrunn svaret.
\hint{Hva må gjelde for $f(1)$, $f(2)$ og $f(4)$?}

\oppg{U6}{3}{\mh}
Kleibers lov (1932) bygger på målinger av mange pattedyr. Den sier at energiforbruket i hvile er omtrent proporsjonalt med $m^{\tfrac34}$, der $m$ er massen.
\begin{tasks}(1)
  \task Hvor mange ganger mer energi bruker et dyr som er 16 ganger så tungt?
  \task Hvor mye energi bruker det store dyret per kilo, sammenlignet med det lille?
  \task Et dyr er 20\,000 ganger så tungt som et annet. Finn forholdet mellom energiforbrukene. Svar med tre gjeldende siffer.
  \task Drøft: Hva begrenser modellen? Er den gyldig for alle dyr?
\end{tasks}

\begin{oppsummering}
En eksponentialfunksjon ganges med en fast faktor når $x$ øker med et fast tall.
En potensfunksjon ganges med en fast faktor når $x$ ganges med et fast tall.
Bruk forhold mellom verdier til å skille typene, og husk at en modell bare er en tilnærming.
\end{oppsummering}
